\subsection{Success and failure cases}\label{sec:cases}
\owner{each owner: one figure per case, with the cause named in the caption}

% Rubric: "illustrate and discuss the successful cases and failure cases", "good visualizations".

\subsubsection{Successes}
\paragraph{HDBSCAN on D2: belts of varying density.} Shown in Figure~\ref{fig:density}. HDBSCAN keeps
the sparse belts of the Americas (62\% of earthquakes clustered, against 20\% for DBSCAN) as long
clusters that follow the plate boundary. It succeeds because its density threshold adapts to each
cluster, the property D2 needs most.
\habegin
% Source: hoanganh/comparison/results/report_d6_d7.csv; figure from hoanganh/comparison/run_report_tables.py.
\paragraph{GMM on D7: classes of unequal spread.} \hatag{} Shown in Figure~\ref{fig:bc_cases}. GMM matches the
diagnosis best of all methods (ARI 0.774). K-means places its boundary too far into the wide malignant
class (ARI 0.654); the full covariance of each GMM component lets the boundary follow the compact benign
class. Cause of success: the model allows clusters of different spread.
\haend
\todo{Lam}{DBSCAN / Spectral on D1.}
\todo{Ky}{K-Means++ on D4: compact classes in 7 dimensions.}
% Sources: weiyew/sc4020-hdb/results/ablation_features.csv, ablation_flat_design_clusters.csv
% (variant "floor area and lease only").
\paragraph{HDBSCAN on D8, floor area and lease only: flat designs.} This feature set was chosen after
seeing the main results, so we present it as an ablation finding and a sanity check, not as a tuned
setting. Dropped to two features, D8 does contain dense spots: standard flat designs, built to fixed
plans in particular years. HDBSCAN finds 14 clusters (43.2\% noise), and 13 of them are at least 97\% one
flat type (Figure~\ref{fig:flat_designs}); in the median cluster one flat model accounts for
83\% of sales. A simple grouping by flat
model and lease would give similar groups, which is the point: the method recovers a structure we can
name. ARI against flat type is still only 0.127, because each type is split into several designs and
the noise counts as a cluster. A low ARI here reflects a finer partition than the labels, not a failure.

\begin{figure}[h]\centering
\includegraphics[width=0.8\textwidth]{hdb_success_flat_designs}
\caption{D8, success of HDBSCAN on floor area and remaining lease only (\texttt{min\_cluster\_size}~200;
grey: noise). Each cluster is one standard flat design; the legend gives its main flat type and model.
Cause of success: fixed building plans create real dense spots in these two features.}
\label{fig:flat_designs}
\end{figure}

\subsubsection{Failures}
\paragraph{K-Means++ on D2: unrelated belts merged.} Shown in Figure~\ref{fig:kmeans_belts}. One
cluster joins five belts whose centres are 2{,}932 to 11{,}449\,km apart, from Iceland to the
Mid-Indian Ridge. Cause: cluster shape. Compact regions cannot follow long belts.
\paragraph{Density methods on D8: one mass, or storey bands.} Shown in Figures~\ref{fig:hdb_cmp}
and~\ref{fig:abl_ties}. At the chosen settings DBSCAN puts 97.1\% of sales in one cluster. Under a cap on
the largest cluster, the clusters become single storey bands. Causes: no gaps between groups, and values
recorded in steps.
\habegin
\paragraph{DBSCAN on D7: half the data discarded.} \hatag{} Shown in Figure~\ref{fig:bc_cases}. The least noisy
two-cluster setting keeps the dense benign core as one cluster, finds a second cluster of 14 malignant tumours and
labels 55.5\% of tumours noise, including 183 of the 212 malignant ones. Its silhouette (0.356) is higher than that
of GMM or K-means, and its ARI is 0.222. Causes: 30 dimensions, and one class much sparser than the other.

\begin{figure}[h]\centering
\includegraphics[width=\textwidth]{bc_cases}
\caption{\hatag D7 on its first two principal components (grey: noise; each cluster takes the colour
of its majority class). Success of GMM against K-means: the boundary follows the compact benign class.
Failure of DBSCAN: only the dense benign core is clustered. Causes: classes of unequal spread, and no
density contrast in 30 dimensions.}
\label{fig:bc_cases}
\end{figure}
\haend
\todo{Lam}{DBSCAN on D5.}
\todo{Ky}{HDBSCAN on D4: 37\% noise removes rare classes.}

\begin{figure}[h]\centering
\figtodo{All}{success/failure panel grid: one panel per case above}
\caption{\todo{All}{caption naming the cause of each case}}
\label{fig:cases}
\end{figure}
